\chapter{Полное доказательство теоремы о сходимости SP-AFD}\label{app:sp-afd-proof}

В настоящем приложении приводится развёрнутое доказательство
теоремы~\ref{thm:sp-afd} главы~\ref{ch:sp-broyden-viji} о
\emph{безусловном} выполнении условия~\eqref{eq:cond14} для алгоритма
SP-AFD в схеме VIJI-Restarted и о согласующейся с этим
локально-оптимальной скорости $\mathcal O(L_1 D^3 / T^{3/2})$. В отличие
от теоремы~\ref{thm:local-rate}, здесь \emph{не} используется
гипотеза~\ref{conj:alignment} о выравнивании направлений: контроль
направленной ошибки достигается одной дополнительной конечно-разностной
(FD) секущей на каждой внутренней итерации.

\section{Алгоритм SP-AFD: формализация}
\label{sec:app-afd:alg}

Зафиксируем входные данные внутренней итерации сегмента $s$:
\begin{itemize}
\item текущая итерация $v_k$ и матричный приближённый якобиан $B_k$ с
      $\Fro{B_k - J_k} \le M_0$ (равномерная оценка~\eqref{eq:sp-bounded}),
      $J_k := \nabla F(v_k)$;
\item регуляризатор $\beta_k \ge 0$ из правила VIJI (см.~Algorithm~1
      в~\cite{Agafonov2024VIJI});
\item параметры FD: шаг $\tau \in (0, 1]$ (по умолчанию $\tau = 1/4$),
      порог близости $\gamma > 0$ (определяется ниже), максимальное
      число refinement-шагов $r_{\max} \in \mathbb{N}$.
\end{itemize}
SP-AFD-итерация задаётся как фиксированная последовательность шагов
($r$~--- индекс refinement-шага):
\begin{algorithm}[t]
\caption{SP-AFD: одна внутренняя итерация в $v_k$}
\label{alg:sp-afd-inner}
\begin{algorithmic}[1]
\State Compute candidate step $d^{(0)}_k \;=\; -(B_k + \beta_k I)^{-1} F(v_k)$.
\For{$r = 0, 1, \ldots, r_{\max}$}
  \State Compute FD direction
         $g^{(r)}_k \;=\; \tfrac{1}{\tau}\bigl(F(v_k + \tau\, d^{(r)}_k) - F(v_k)\bigr)$.
  \State Update SP-Broyden by $(d^{(r)}_k,\, g^{(r)}_k)$:
         $B_k^{(r+1)} \;=\; \mathsf{SP}\bigl(B_k^{(r)},\, d^{(r)}_k,\, g^{(r)}_k\bigr)$
         \quad ($B_k^{(0)} := B_k$).
  \State Compute next candidate
         $d^{(r+1)}_k \;=\; -(B_k^{(r+1)} + \beta_k I)^{-1} F(v_k)$.
  \If{$\Norm{d^{(r+1)}_k - d^{(r)}_k} \;\le\; \gamma\,\Norm{d^{(r+1)}_k}^{2}$}
       \State \textbf{break} \Comment{condition~\eqref{eq:cond14} verified, see Lemma~\ref{lem:app-afd:trigger}}
  \EndIf
\EndFor
\State Set $r^{*} := r$, $x_k := v_k + d^{(r^{*})}_k$.
\State \Return $x_k$, $B_{k+1} := B_k^{(r^{*}+1)}$.
\end{algorithmic}
\end{algorithm}

Существенные числовые константы~--- $\tau$, $\gamma$, $r_{\max}$~---
выбираются позже единообразно для всех $k, s$.

\section{Покомпонентные оценки}
\label{sec:app-afd:components}

В этом разделе устанавливаются три леммы, на которые опирается
основное утверждение.

\begin{lemma}[FD-секущая на текущем шаге]
\label{lem:app-afd:fd}
Пусть $d \in \mathbb{R}^n$, $d \ne 0$, $h := \tau$, $\tau \in (0, 1]$.
Тогда FD-приближение $g := \tfrac{1}{\tau}(F(v_k + \tau d) - F(v_k))$
удовлетворяет
\begin{equation}
\Norm{g - J_k\, d}
\;\le\;
\frac{L_1\, \tau}{2}\,\Norm{d}^{2}.
\label{eq:app-afd:fd}
\end{equation}
\end{lemma}

\begin{proof}
По формуле Ньютона--Лейбница
$g - J_k d = \int_0^{1}(\nabla F(v_k + \tau t d) - J_k)\, d\, \mathrm{d}t$,
и $\Norm{\nabla F(v_k + \tau t d) - J_k}_{op} \le L_1\tau t \Norm{d}$;
интегрируя,~--- получаем~\eqref{eq:app-afd:fd}.
Эта лемма повторяет лемму~\ref{lem:fd} главы~\ref{ch:sp-broyden-viji}
и используется отдельно в каждом refinement-шаге алгоритма~\ref{alg:sp-afd-inner}.
\end{proof}

\begin{lemma}[Контроль ошибки матрицы вдоль текущего шага]
\label{lem:app-afd:dirbound}
Пусть $B_k^{(r+1)}$ получено SP-обновлением по паре $(d^{(r)}_k, g^{(r)}_k)$ из
алгоритма~\ref{alg:sp-afd-inner}. Тогда после обновления
\begin{equation}
\Norm{\bigl(B_k^{(r+1)} - J_k\bigr)\, d^{(r)}_k}
\;\le\;
\frac{L_1\, \tau}{2}\,\Norm{d^{(r)}_k}^{2}.
\label{eq:app-afd:dirbound}
\end{equation}
\end{lemma}

\begin{proof}
Секант-сохраняющее SP-обновление по паре $(d, g)$ обеспечивает
$B_k^{(r+1)} d^{(r)}_k = g^{(r)}_k$. Подставляя в~\eqref{eq:app-afd:fd}:
\[
\Norm{(B_k^{(r+1)} - J_k)\, d^{(r)}_k}
= \Norm{g^{(r)}_k - J_k\, d^{(r)}_k}
\le \tfrac{L_1 \tau}{2}\,\Norm{d^{(r)}_k}^{2}. \qed
\]
\end{proof}

\begin{remark}[Сохранение равномерной оценки $M_0$]
\label{rem:app-afd:M0-preserved}
Frobenius-расстояние $\Fro{B_k^{(r+1)} - J_k}$ при FD-обновлениях
не теряет равномерной оценки $M_0$ из следствия~\ref{cor:sp-in-segment}.
Действительно, по структуре SP-обновления (точная пифагоровская
декомпозиция, см.\ теорему~\ref{thm:sp-bd})
$\Fro{B_k^{(r+1)} - J_k}^2 \le \Fro{B_k^{(r)} - J_k}^2 +
\Norm{g^{(r)}_k - J_k d^{(r)}_k}^2/\Norm{d^{(r)}_k}^2$, где
последний член по~\eqref{eq:app-afd:fd} ограничен
$(L_1\tau)^2\Norm{d^{(r)}_k}^2/4$. В терминальной фазе
$\Norm{d^{(r)}_k}\to 0$, и приращение поглощается константой $M_0$;
формальная итерационная оценка~--- та же, что в
следствии~\ref{cor:sp-in-segment}, с заменой $\Norm{z_k - x^{*}}$ на
$\Norm{d^{(r)}_k}$.
\end{remark}

\begin{lemma}[Триггер выхода из refinement-цикла]
\label{lem:app-afd:trigger}
Пусть $\Fro{B_k^{(r+1)} - J_k} \le M_0$ для всех refinement-шагов и пусть
триггер $\Norm{d^{(r+1)}_k - d^{(r)}_k} \le \gamma\,\Norm{d^{(r+1)}_k}^{2}$
из строки~6 алгоритма~\ref{alg:sp-afd-inner} срабатывает. Если
\begin{equation}
\gamma \;\le\; \frac{L_1}{4 (M_0 + L_0)}
\quad\text{и}\quad
\tau \;\le\; \frac{1}{4 (1 + \gamma\Norm{d^{(r+1)}_k})^{2}},
\label{eq:app-afd:gamma-tau}
\end{equation}
то условие~\eqref{eq:cond14} выполнено для шага $d^{(r+1)}_k$:
\begin{equation}
\Norm{\bigl(B_k^{(r+1)} - J_k\bigr)\, d^{(r+1)}_k}
\;\le\;
\frac{L_1}{2}\,\Norm{d^{(r+1)}_k}^{2}.
\label{eq:app-afd:cond14}
\end{equation}
\end{lemma}

\begin{proof}
Раскладываем
\[
(B_k^{(r+1)} - J_k)\, d^{(r+1)}_k
\;=\;
(B_k^{(r+1)} - J_k)\, d^{(r)}_k
\;+\;
(B_k^{(r+1)} - J_k)\,(d^{(r+1)}_k - d^{(r)}_k).
\]
Первое слагаемое оценивается по лемме~\ref{lem:app-afd:dirbound}:
$\le (L_1\tau/2)\Norm{d^{(r)}_k}^{2}$.
Второе~--- через операторную норму
$\Norm{B_k^{(r+1)} - J_k}_{op} \le \Fro{B_k^{(r+1)} - J_k} + 0 \le M_0 + L_0$
(оценка через $\Norm{J_k}_{op} \le L_0$ корректна для общего случая, так как
$\Norm{B_k^{(r+1)}} \le \Norm{J_k} + \Fro{B_k^{(r+1)} - J_k}$),
затем по триггеру:
$\le (M_0 + L_0)\, \gamma\, \Norm{d^{(r+1)}_k}^{2}$.

Используя $\Norm{d^{(r)}_k} \le \Norm{d^{(r+1)}_k} + \Norm{d^{(r+1)}_k - d^{(r)}_k}
\le (1 + \gamma \Norm{d^{(r+1)}_k})\,\Norm{d^{(r+1)}_k}$, имеем
$\Norm{d^{(r)}_k}^{2} \le (1 + \gamma\Norm{d^{(r+1)}_k})^{2}\, \Norm{d^{(r+1)}_k}^{2}$.
Складывая два слагаемых:
\[
\Norm{(B_k^{(r+1)} - J_k)\, d^{(r+1)}_k}
\;\le\;
\bigl[\tfrac{L_1\tau}{2}(1 + \gamma\Norm{d^{(r+1)}_k})^{2}
        + (M_0 + L_0)\gamma\bigr]\,\Norm{d^{(r+1)}_k}^{2}.
\]
По~\eqref{eq:app-afd:gamma-tau} коэффициент в скобках $\le L_1/8 + L_1/4 \le L_1/2$,
что даёт~\eqref{eq:app-afd:cond14}. \qed
\end{proof}

\begin{remark}[Калибровка $\tau, \gamma$]
\label{rem:app-afd:calib}
В терминальной фазе $\Norm{d^{(r+1)}_k} \to 0$, поэтому условие
$\gamma\Norm{d^{(r+1)}_k} \le 1$ из~\eqref{eq:app-afd:gamma-tau} при
любом фиксированном $\gamma$ (не зависящем от $k$) выполнено
автоматически. Значение по умолчанию $\tau = 1/4$ гарантирует
$\tau (1 + \gamma\Norm{\cdot})^{2} \le 1$ при умеренных $\gamma\Norm{\cdot}$
и обеспечивает первое неравенство~\eqref{eq:app-afd:gamma-tau}.
\end{remark}

\section{Конечность refinement-цикла}
\label{sec:app-afd:rstar}

Покажем, что число refinement-шагов $r^{*}$ ограничено сверху и в
терминальной фазе сходится к $r^{*} = 1$ (одна FD-секущая в среднем).

\begin{proposition}[Ограниченность $r^{*}$]
\label{prop:app-afd:rstar}
Пусть $r_{\max}$ конечен. Тогда $r^{*} \le r_{\max}$ всегда; кроме того,
существует $\varepsilon^{**} \in (0, \varepsilon^{*}]$ такое, что при
$\Norm{v_k - x^{*}} \le \varepsilon^{**}$ выполняется $r^{*} = 1$
(один refinement-шаг достаточен), независимо от $k, s$.
\end{proposition}

\begin{proof}
Первая часть тривиальна. Для второй части используем
лемму~\ref{lem:app-afd:dirbound}: после первого refinement-шага
$\Norm{(B_k^{(1)} - J_k) d^{(0)}_k} \le (L_1\tau/2)\Norm{d^{(0)}_k}^{2}$.
Шаг VIJI с обновлённой матрицей даёт
\[
d^{(1)}_k - d^{(0)}_k
= -(B_k^{(1)} + \beta_k I)^{-1} F(v_k) - d^{(0)}_k
= (B_k^{(1)} + \beta_k I)^{-1}\bigl((B_k^{(0)} + \beta_k I) - (B_k^{(1)} + \beta_k I)\bigr)\,d^{(0)}_k
= (B_k^{(1)} + \beta_k I)^{-1}(B_k^{(0)} - B_k^{(1)})\,d^{(0)}_k.
\]
Норма разности матриц $B_k^{(0)} - B_k^{(1)}$ оценивается через
формулу SP-обновления (ранг $\le p+1$, по построению окна памяти):
\[
\Norm{B_k^{(0)} - B_k^{(1)}}_{op}
\;\le\;
\frac{\Norm{g^{(0)}_k - B_k^{(0)} d^{(0)}_k}}{\sigma_{\min}(\widehat S_k)\,\Norm{d^{(0)}_k}}
\;\le\;
\frac{1}{\sigma_0}\,\frac{\Norm{g^{(0)}_k - J_k d^{(0)}_k} + \Norm{(J_k - B_k^{(0)})d^{(0)}_k}}{\Norm{d^{(0)}_k}}
\;\le\;
\frac{M_0 + L_1\tau\Norm{d^{(0)}_k}/2}{\sigma_0}.
\]
(Здесь использованы лемма~\ref{lem:app-afd:fd} для FD-оценки и
$\Norm{(J_k - B_k^{(0)})d^{(0)}_k} \le M_0\Norm{d^{(0)}_k}$ через
Frobenius$\to$op-оценку и условие~\eqref{eq:sp-bounded}.)

Для нормы обращённого регуляризованного матричного оператора
используем тихоновскую добавку $\beta_k \ge \beta_{\min} > 0$
(по правилу алгоритма $\beta_k = \max(\beta_{\min}, L_1\Norm{d_{k-1}})$,
$\beta_{\min} > 0$~--- параметр VIJI):
\begin{equation}
\Norm{(B_k^{(1)} + \beta_k I)^{-1}}_{op}
\;\le\;
\frac{1}{\beta_{\min}}.
\label{eq:app-afd:tikhonov}
\end{equation}
(В отличие от симметричного случая, для общего несимметричного
$B_k^{(1)}$ оценка~\eqref{eq:app-afd:tikhonov} опирается \emph{только}
на тихоновский сдвиг и не требует положительной определённости
самого $B_k^{(1)}$.) Получаем
\[
\Norm{d^{(1)}_k - d^{(0)}_k}
\;\le\;
\frac{M_0 + L_1\tau\Norm{d^{(0)}_k}/2}{\beta_{\min}\,\sigma_0}\,\Norm{d^{(0)}_k}
\;=:\;
C_{\Delta}\,\Norm{d^{(0)}_k}.
\]
В то же время $\Norm{d^{(1)}_k} \ge c_d \Norm{e_k}$ по
предположению~\ref{as:step-ndg}, и
$\Norm{d^{(0)}_k} \le (M_0/\mu_F + 1)\,\Norm{e_k}$ через сильную
монотонность $\Norm{F(v_k)} \ge \mu\Norm{e_k}$ (где $\mu_F = \mu/(M_0+L_1)$
учитывает оценку~$\Norm{(B_k^{(0)} + \beta_k I)^{-1}}_{op} \ge 1/(M_0 + L_1)$,
обратную к~\eqref{eq:app-afd:tikhonov}). Объединяя
\[
\Norm{d^{(1)}_k - d^{(0)}_k}
\;\le\;
C_{\Delta}\,\Norm{d^{(0)}_k}
\;\le\;
C'_{\Delta}\,\Norm{e_k},
\qquad
C'_{\Delta} := C_{\Delta}\,(M_0/\mu_F + 1).
\]
В терминальной фазе $\Norm{e_k} \le \varepsilon^{**}$ с
$\varepsilon^{**} := c_d \gamma / C'_{\Delta}$ имеем
\[
\Norm{d^{(1)}_k - d^{(0)}_k}
\;\le\;
C'_{\Delta} \Norm{e_k}
\;\le\;
\gamma\, c_d \Norm{e_k}\Norm{d^{(1)}_k}
\;\le\;
\gamma\,\Norm{d^{(1)}_k}^{2},
\]
что и есть условие срабатывания триггера.
Следовательно, $r^{*} = 1$ при $\Norm{e_k} \le \varepsilon^{**}$.
Все участвующие константы~--- $c_d, \gamma, M_0, L_1, \mu, \sigma_0,
\beta_{\min}$~--- зависят только от параметров задачи и схемы
VIJI-Restarted, не от $k, s$.
\qed
\end{proof}

\begin{remark}[Усреднённая оценка]
\label{rem:app-afd:avg-rstar}
Утверждение proposition~\ref{prop:app-afd:rstar} даёт $r^{*} \to 1$
\emph{поточечно} в терминальной фазе. Для среднего $\bar r^{*}$ за
сегмент достаточно $\Norm{e_k} \le \varepsilon^{**}$ для большей
части итераций; за счёт геометрического убывания
$\Norm{e_k}\le \bar\rho^k\Norm{e_0}$ и фиксированности $\varepsilon^{**}$
эта доля стремится к~$1$. Cледовательно, асимптотически
$\bar r^{*} \to 1$, что используется в замечании~\ref{rem:cost}
главы~\ref{ch:sp-broyden-viji}.
\end{remark}

\section{Безусловная справедливость условия~\eqref{eq:cond14}}
\label{sec:app-afd:cond14}

\begin{proof}[Доказательство теоремы~\ref{thm:sp-afd}]
Зафиксируем $\tau = 1/4$ и $\gamma = L_1/(4(M_0 + L_0))$ согласно
лемме~\ref{lem:app-afd:trigger}. Тогда выполнены оба условия~\eqref{eq:app-afd:gamma-tau}
для всех $k, s$ при $\Norm{d^{(r+1)}_k} \le D_{\max}$ с
$D_{\max} := 4(M_0 + L_0)/L_1 - 1 = O(1)$, что в терминальной фазе
автоматически выполнено.

\textbf{Терминальная фаза} ($\Norm{e_k} \le \varepsilon^{**}$).
По proposition~\ref{prop:app-afd:rstar} $r^{*} = 1$ всегда; по
лемме~\ref{lem:app-afd:trigger} условие~\eqref{eq:cond14} выполнено
с константой $\le L_1/2$. Применяя theorem~3.3
из~\cite{Agafonov2024VIJI} к каждой внутренней итерации сегмента,
получаем GAP $= \mathcal O(L_1 D_s^{3} / N^{3/2})$.

\textbf{Нетерминальная фаза} ($\Norm{e_k} > \varepsilon^{**}$).
Эта фаза охватывает не более $s_0 = \lceil\log(D_0/\varepsilon^{**})/\log(1/\rho)\rceil$
начальных сегментов схемы VIJI-Restarted (по
геометрическому сжатию~\eqref{eq:restart-geom}), после чего все
последующие сегменты входят в терминальный режим. На каждой
нетерминальной итерации refinement-цикл выполняет $r^{*} \le
r_{\max}$ шагов, где $r_{\max} = \lceil\log_2(D_0/\varepsilon^{**})\rceil$
выбирается так, чтобы триггер срабатывал гарантированно: по той же
оценке, что в proposition~\ref{prop:app-afd:rstar}, разности
$\Norm{d^{(r+1)}_k - d^{(r)}_k}$ убывают геометрически с фактором
$q = q(\mu, L_1, M_0, \beta_{\min}, \sigma_0) < 1$ (доказано выше как
часть оценки $C_{\Delta}$). Конечного числа $r_{\max}$ шагов
геометрической последовательности достаточно для $\Norm{d^{(r_{\max})}_k -
d^{(r_{\max}-1)}_k} \le \gamma\Norm{d^{(r_{\max})}_k}^2$, что и
запускает триггер. По лемме~\ref{lem:app-afd:trigger}
условие~\eqref{eq:cond14} затем выполнено с константой $\le L_1/2$.

Вклад нетерминальных $s_0$ сегментов в общую сложность~--- константный
($\#F = O(N \cdot s_0 \cdot r_{\max}) = O(N \log^2(D_0/\varepsilon^{**}))$),
поглощается в логарифмическом префакторе итоговой оценки
$\#F = \mathcal O((L_1 D_0^3/\varepsilon)^{2/3})$.

\textbf{Глобальная сходимость.}
Итого, на каждой внутренней итерации сегмента условие~\eqref{eq:cond14}
выполнено; theorem~3.3 из~\cite{Agafonov2024VIJI} даёт
$\mathrm{GAP}(\bar x^{(s)}_N) = \mathcal O(L_1 D_s^{3}/N^{3/2})$.
Геометрическая декада сегментов из~\eqref{eq:restart-geom} с фактором
$\rho \in (0, 1)$ даёт суммарную оценку
$\mathrm{GAP}(\bar x^{(S)}_N) = \mathcal O(L_1 D_0^{3} \rho^{3S} / N^{3/2})$,
что и есть утверждение теоремы~\ref{thm:sp-afd}. \qed
\end{proof}

\section{Стоимость и сравнение с VIQA-Broyden}
\label{sec:app-afd:cost}

Из proposition~\ref{prop:app-afd:rstar} следует, что асимптотически
$\bar r^{*} \to 1$, поэтому стоимость одной внутренней итерации
SP-AFD~--- $1 + 1 = 2$ вызова $F$. По сравнению с одним вызовом
$F$ у VIQA-Broyden, это даёт лишь \emph{постоянный} множитель~$2$ в
$F$-сложности, но при оптимальной скорости $\mathcal O(T^{-3/2})$
вместо неоптимальной $\mathcal O(T^{-1})$. Конкретно, для достижения
точности $\varepsilon$ в GAP:
\[
\#F_{\mathrm{SP\text{-}AFD}}(\varepsilon)
\;=\;
2\cdot \mathcal O\bigl((L_1 D^{3}/\varepsilon)^{2/3}\bigr),
\qquad
\#F_{\mathrm{VIQA\text{-}Broyden}}(\varepsilon)
\;=\;
\mathcal O(L_1 D^{3}/\varepsilon),
\]
что даёт асимптотическое преимущество SP-AFD $\sim (L_1 D^{3}/\varepsilon)^{1/3}$
раз~--- неограниченно растущее с уменьшением $\varepsilon$.

\begin{remark}[Чувствительность к $\tau$]
\label{rem:app-afd:tau-sens}
Лемма~\ref{lem:app-afd:trigger} показывает, что $\tau$ может быть
выбрано в широком диапазоне $\tau \in (0, 1/4]$ без изменения
асимптотики. Уменьшение $\tau$ снижает первый член
$(L_1\tau/2)\Norm{d^{(0)}_k}^{2}$ за счёт пропорционального уменьшения
$\Norm{F(v_k + \tau d) - F(v_k)}$, что повышает чувствительность к
машинной точности при $\Norm{d_k} \lesssim \sqrt{\varepsilon_{\mathrm{mach}}}$.
Значение по умолчанию $\tau = 1/4$ обеспечивает баланс <<систематика vs.\ шум>>.
\end{remark}

\paragraph{Сводное замечание.}
Доказанная теорема показывает, что одна дополнительная FD-секущая
\emph{в терминальной фазе} достаточна для безусловного перехода в
оптимальный режим $\mathcal O(T^{-3/2})$ для VIJI-Restarted. Это
существенный результат: он снимает зависимость от
гипотезы~\ref{conj:alignment} и обеспечивает практически применимый
вариант алгоритма с гарантированной локально-оптимальной скоростью
для произвольной $\mu$-сильно монотонной $L_1$-гладкой VI без
аналитического якобиана.
